\cotefile{15}
\scanpage{49}
\mspage{21}
% archive sheet number: 48
Tout revient donc à prouver qu'il existe $q$ séries entières absolument sommables par rapport aux $t_i$ et sans terme constant, qui satisfont aux équations
\[
v_k=\frac{\unsure{M|\mathbb{M}}}{(1-\handwritten{\mu}\handwritten{\sum}_{\handwritten{j=1}}^{p}t_j)(1-\handwritten{\mu}\handwritten{\sum}_{\handwritten{i=1}}^{q}v_i)}-M-M\handwritten{\mu}\handwritten{\sum}_{\handwritten{i=1}}^{q}v_i
\qquad (1\handwritten{\leq}k\handwritten{\leq}q)
\tag*{(9)}
\]

Or, on déduit aussitôt \deleted{que}\inserted{\handwritten{de}} (9) que si ce système a une solution,tous les $v_i$ sont égaux à $v_1=v$ sa tisfaisant à
\[
(1-q\handwritten{\mu}v)(1+qM\handwritten{\mu})v=\frac{M}{1-\handwritten{\mu}\handwritten{\sum}_{\handwritten{j=1}}^{p}t_j}-M(1-q\handwritten{\mu}v)
\tag*{(10)}
\]
et réciproquement ; or, la fonction $v(t_1,\ldotp\ldotp,t_p)$ qui s'annule pour les $t_j$ nuls et satisfait à (10) est donnée par
\[
v(t_1,\ldotp\ldotp,t_p)=\frac{1}{2q\handwritten{\mu}(1+qM\handwritten{\mu})}\handwritten{\Biggl[}1-\sqrt{1-4Mq\handwritten{\mu}(1+qM\handwritten{\mu})\left(\frac{1}{1-\handwritten{\mu}\handwritten{\sum}_{\handwritten{j=1}}^{p}t_j}-1\right)}\handwritten{\Biggr]}
\]
Comme la fonction $(1-w)^{\handwritten{\frac{1}{2}}}$ est égale à la somme de la série entière $\handwritten{\sum}_{\handwritten{n=0}}^{\handwritten{\infty}}\binom{\handwritten{\frac{1}{2}}}{n}v^n$ pour $-1\handwritten{<}w\handwritten{<}1$, il résulte du th. de substitution (§1,th.2) que $v(t_1,\ldotp\ldotp,t_p)$ est égale à une série entière absolument sommable et sans terme constant dans un polycylindre de centre $(0,\ldots,0)$.

Reste à montrer que si $w_i$ $(1\handwritten{\leq}i\handwritten{\leq}q)$ sont $q$ fonctions \uline{continues} au voisinage de l'origine dans ${\unsure{K|\mathbb{K}}}^{p}$, nulles en ce point et satisfaisant aux relations
\[
w_i=h_{i0}+\handwritten{\sum}_{\handwritten{(n_j)}}h_{i;n_1\ldotp\ldotp n_q}w_1^{n_1}\ldots w_q^{n_q}
\qquad\text{pour }1\handwritten{\leq}i\handwritten{\leq}q
\tag*{(11)}
\]
dans un voisinage de l'origine dans $K^{p}$, on a nécessairement $u_i\unsure{\leq|=}w_i$ da ns\handwritten{\unsure{.}} un voisinage de l'origine pour $1\handwritten{\leq}i\handwritten{\leq}q$. Or, compte-tenu de la prop.8 du §1, on tire de (5) et (11) que
\[
u_i-w_i=\underline{o}\bigl(\sup_{1\handwritten{\leq}k\handwritten{\leq}q}\handwritten{|}u_k-\unsure{w|v}_k\handwritten{|}\bigr)
\qquad\text{pour }1\handwritten{\leq}i\handwritten{\leq}q
\]
dans un voisinage de $0$, ce qui entraîne $u_i=w_i$ pour $1\handwritten{\leq}i\handwritten{\leq}q$ dans un voisinage de $0$ (\uline{Fonct.var.réelle}, chap.V,§2). Le th.1 est ainsi complètement démontré.

\scanpage{51}
\mspage{13}
% archive sheet number: 50
\[
\begin{aligned}
\handwritten{f}(u_1+w_1,\ldotp\ldotp,u_p+w_p)={}&\handwritten{f}(u_1,\ldotp\ldotp,u_p)+\handwritten{\sum}_{\handwritten{k=1}}^{p}D_k\handwritten{f}(0,\ldotp\ldotp,0)w_k+\handwritten{\sum}_{\handwritten{k=1}}^{p}\handwritten{f}_k(u_1,\ldotp\ldotp,u_p)w_k+\\
&+\handwritten{\sum}_{\handwritten{(n_j)(m_j)}}\handwritten{b}(n_1,\ldotp\ldotp,n_p,m_1,\ldotp\ldotp,m_p)u_1^{n_1}\ldotp\ldotp u_p^{n_p}w_1^{m_1}:\ldotp\ldotp w_p^{m_p}
\end{aligned}
\]
où les séries entières (absolument sommables) $\handwritten{f}_k$ n'ont pas de terme constant et où, dans la dernière série (absolument sommable) on a $\handwritten{\sum}_{\handwritten{k=1}}^{p}m_k\handwritten{\geq}2$ pour tous les termes. Supposons les séries absolument sommables pour $\handwritten{|}u_k\handwritten{|}=\handwritten{|}w_k\handwritten{|}=r$ $(1\handwritten{\leq}k\handwritten{\leq}p)$, et soit $\unsure{M|\mathbb{M}}$ la borne supérieure des termes $\handwritten{\|}\handwritten{b}(n_1,\ldotp\ldotp,n_p,m_1,\ldotp\ldotp,m_p)\handwritten{\|}r^{\handwritten{\sum}_{k}(n_k+m_k)}$. Alors pour $\handwritten{|}u_k\handwritten{|}\handwritten{\leq}\handwritten{\frac{1}{2}}r$ et $\handwritten{|}w_k\handwritten{|}\handwritten{\leq}\handwritten{\frac{1}{2}}r$ $(1\handwritten{\leq}k\handwritten{\leq}p)$ on a, en posant $t=\sup_{\underset{\handwritten{\unsure{\smile}}}{k}}\handwritten{|}w_k\handwritten{|}$
\[
\begin{aligned}
&\handwritten{\sum}_{\handwritten{(n_j)(m_j)}}\handwritten{\|}\handwritten{b}(n_1,\ldotp\ldotp,n_p,m_1,\ldotp\ldotp,m_p)\handwritten{\|}\handwritten{\bigl|}u_1^{n_1}\ldots u_p^{n_p}w_1^{m_1}\ldotp\ldotp w_p^{m_p}\handwritten{\bigr|}\handwritten{\leq}\\
&\handwritten{\leq}Mt^2\handwritten{\sum}_{\handwritten{(n_j)(m_j)}}\left(\handwritten{\frac{1}{2}}\right)^{\handwritten{\sum}_{k}n_k}\left(\handwritten{\frac{t}{r}}\right)^{(\handwritten{\sum}_{k}m_k)-2}
\end{aligned}
\]
Comme la série qui figure au second membre de cette inégalité est convergente pour $\handwritten{|}t\handwritten{|}\handwritten{\leq}\handwritten{\frac{1}{2}}r$, on voit que la formule (6) résulte aussitôt de la continuité des séries entières.

COROLLAIRE.- \uline{Toute fonction holomorphe dans} $A$ \uline{est indéfiniment différentiable dans} $A$.

En effet, la prop.8 montre que toute fonction holomorphe $\handwritten{f}$ est différentiable en tout point, et que sa différentielle est l'application
\[
(u_1,\ldots,u_p)\handwritten{\longrightarrow}\handwritten{\sum}_{\handwritten{k=1}}^{p}D_k\handwritten{f}\,.u_k
\]
qu'on peut identifier à l'élément $(D_k\handwritten{f}(z_1,\ldotp\ldotp,z_p))_{1\handwritten{\leq}k\handwritten{\leq}p}$ de l'\uline{espace produit} ${\unsure{E|\mathbb{E}}}^{p}$ en tout point. Or, il résulte de la prop.7 que cet élément est fonction holomorphe dans $A$, ce qui prouve que $\handwritten{f}$ est indéfiniment différentiable dans $A$.

Pour $p=1$, on dit qu'un point $a\handwritten{\in}\unsure{K|\mathbb{K}}$ est un \uline{zéro d'ordre $h$} d'une fonction $\handwritten{f}$ holomorphe au voisinage de $a$ si, dans le développement de $\handwritten{f}$ en série entière suivant les puissances de $z-a$, le premier terme non nul a un exposant égal à $h$. Si $K$ est de caractéristique $0$, il est

\scanpage{53}
\mspage{12}
% archive sheet number: 52
En effet, en désignant par $\handwritten{\unsure{c|C}}_{k;n_1,\ldotp\ldotp,n_p}$ les coeffici\handwritten{\`}{e}nts de $\handwritten{f}_k$, les coefficients $\handwritten{b}_{n_1\ldotp\ldotp n_p}$ de $\handwritten{g}$ (supposée exister) doivent satisfaire aux conditions $n_k\handwritten{b}_{n_1\ldotp\ldotp n_p}=\handwritten{c}_{k;n_1,\ldotp\ldotp,n_k-1,\ldotp\ldotp,n_p}$ pour $(n_1,\ldotp\ldotp,n_p)\handwritten{\neq}\handwritten{\neq}(0,\ldotp\ldotp,0)$ et $\handwritten{b}_{00\ldotp\ldotp0}=\handwritten{c}$ ; l'hypothèse implique que ces conditions déterminent les $\handwritten{b}_{n_1\ldotp\ldotp n_p}$ de façon unique. En outre, on a, dans $\unsure{K|\mathbb{K}}$, $\handwritten{\bigl|}\frac{1}{n}\handwritten{\bigr|}\handwritten{\leq}n^{\handwritten{\beta}}$ pour une constante fixe $\handwritten{\beta}>0$ ne dépendant que de la valeur absolue de $K$, et on vérifie aisément que dans $\handwritten{\mathbb{R}}$, la série de terme général,
\[
((n_1+1)(n_2+1)\ldotp\ldotp(n_p+1))^{\handwritten{\beta}}\handwritten{\|}\handwritten{a}_{n_1n_2\ldotp\ldotp n_p}\handwritten{\|}r_1^{n_1}r_2^{n_2}\ldots r_p^{n_p}
\]
converge en même temps que la série de terme général
\[
\handwritten{\|}\handwritten{a}_{n_1\ldotp\ldotp n_p}\handwritten{\|}r_1^{n_1}r_2^{n_2}\ldots r_p^{n_p}.
\]

PROPOSITION 8.- \uline{Soit} $\handwritten{f}(z_1,\ldotp\ldotp,z_p)$ \uline{une fonction holomorphe dans un ensemble ouvert} $A$. \uline{Pour tout point} $\handwritten{a}=(a_1,\ldotp\ldotp,a_p)$, \uline{de} $A$ \uline{et tout nombre} $\handwritten{\varepsilon}>0$, \uline{il existe un polycylindre ouvert} $P_o$ \uline{de centre} $\handwritten{a}$ \uline{tel que}, \uline{pour tout couple de points} $(u_1,\ldotp\ldotp,u_p),(v_1,\ldotp\ldotp,v_p)$ \uline{de ce polycylindre}, \uline{on ait}
\[
\handwritten{\Bigl\|}\handwritten{f}(u_1,\ldotp\ldotp,u_p)-\handwritten{f}(v_1,\ldotp\ldotp,v_p)-\handwritten{\sum}_{\handwritten{k=1}}^{p}D_k\handwritten{f}(a_1,\ldotp\ldotp,a_p)(u_k-v_k)\handwritten{\Bigr\|}\handwritten{\leq}\handwritten{\varepsilon}.\sup_{1\handwritten{\leq}k\handwritten{\leq}p}\handwritten{|}u_k-v_k\handwritten{|}
\tag*{(6)}
\]

On peut se ramener au cas où $\handwritten{a}$ est l'origine, $\handwritten{f}(z_1,\ldotp\ldotp,z_p)$\unsure{'} étant la somme d'une série entière $(\handwritten{c}_{n_1\ldotp\ldotp n_p}z_1^{n_1}\ldotp\ldotp z_p^{n_p})$ absolument sommable dans un polycylindre ouvert $Q$. En vertu du th.2, le polynome $u_k+w_k$ est substituable à $z_k$ dans cette série entière (pour $1\handwritten{\leq}k\handwritten{\leq}p$), dès que $(u_k)$ et $(w_k)$ appartiennent à un polycylindre $P$ assez petit, et par substitution, on obtient une série entière absolument sommable dans $P\handwritten{\times}P$ par rapport à ces $2p$ variables. Par le théorème d'associativité, on peut donc écrire

\scanpage{55}
\mspage{10}
% archive sheet number: 54
\uline{sa somme est une fonction holomorphe dans ce polycylindre}.

De même, en vertu du cor. du th.2 :

PROPOSITION 6.- \uline{Soit} $B$ \uline{un ensemble ouvert dans} ${\unsure{K|\mathbb{K}}}^{q}$, $g_k$ $(1\handwritten{\leq}k\handwritten{\leq}p)$ $p$ \uline{fonctions scalaires holomorphes dans} $B$, $A$ \uline{un ensemble ouvert dans} $K^{p}$ \uline{contenant l'image de} $B$ \uline{par} $(g_1,\ldots,g_p)$. \uline{Si} $\handwritten{f}$ \uline{est une fonction holomorphe dans} $A$, \uline{la fonction composée} $\handwritten{f}(g_1,\ldots,g_p)$ \uline{est holomorphe dans} $B$.

De façon imagée, toute fonction \uline{composée} de fonctions holomorphes est holomorphe.

PROPOSITION 6 bis.- \uline{Soient} $\unsure{E|\mathbb{E}}_1,\ldots,E_r$ \uline{et} $F$ $r+1$ \uline{espaces de Banach sur} $K$, \uline{et soit} $(z_1,\ldotp\ldotp,z_r)\handwritten{\longrightarrow}\handwritten{[}z_1.z_2\ldots z_r\handwritten{]}$ \uline{une application multilinéaire continue de} $E_1\handwritten{\times}\ldots\handwritten{\times}E_r$ \uline{dans} $F$. \uline{Si} $\handwritten{f}_1,\ldots,\handwritten{f}_r$ \uline{sont} $r$ \uline{fonctions holomorphes dans un ensemble ouvert} $A\handwritten{\subset}K^p$, \uline{à valeurs dans} $E_1,\ldots,E_r$ \uline{respectivement}, \uline{la fonction} $\handwritten{[}\handwritten{f}_1.\handwritten{f}_2\ldots\handwritten{f}_r\handwritten{]}$ \uline{est une fonction holomorphe dans} $A$, \uline{à valeurs \handwritten{/}dans} $F$.

Bornons-nous pour simplifier au cas $r=2$. Il suffit alors de remarqu\unsure{er} que si $(x_{\handwritten{\lambda}})_{\handwritten{\lambda\in}L}$ est une famille absolument sommable dans $E_1$ de somme $s$, et $(\handwritten{y}_{\handwritten{\mu}})_{\handwritten{\mu\in}\unsure{M|\mathbb{M}}}$ une famille absolument sommable dans $E_2$, de somme $\handwritten{t}$, la famille $(\handwritten{[}x_{\handwritten{\lambda}}.\handwritten{y}_{\handwritten{\mu}}\handwritten{]})_{(\handwritten{\lambda},\handwritten{\mu})\handwritten{\in}L\handwritten{\times}M}$ est absolument sommable dans $F$ et a pour somme $\handwritten{[}\handwritten{s}.\handwritten{t}\handwritten{]}$, puisqu'il existe.une constante $k\handwritten{>}0$ telle que $\handwritten{\|[}z_{\handwritten{1}}.z_{\handwritten{2}}\handwritten{]\|}\handwritten{\leq}k.\handwritten{\|}z_{\handwritten{1}}\handwritten{\|}.\handwritten{\|}z_{\handwritten{2}}\handwritten{\|}$ ; on appliquera cette remarque aux séries entières dont les sommes sont égales à $\handwritten{f}_1$ et $\handwritten{f}_2$ au voisinage d'un point de $A$.

COROLLAIRE.- \uline{Soit} $\handwritten{f}=(f_1,\ldots,f_q)$ \uline{une fonction définie dans un ensemble ouvert} $A\handwritten{\subset}K^p$ \uline{et à valeurs dans} $K^q$. \uline{Pour que} $\handwritten{f}$ \uline{soit holomorphe dans} $A$, \uline{il faut et il suffit que chacune des fonctions scalaires} $f_k$ $(1\handwritten{\leq}k\handwritten{\leq}q)$ \uline{le soit}.

\scanpage{57}
\mspage{9}
% archive sheet number: 56
 des séries entières $g_k$, et en substituant à $z_k$ dans $\handwritten{f}$ chacune des séries entières en $v_1,\ldotp\ldotp,v_r$ obtenues : $2^{\mathrm{o}}$ soit en substituant $g_k$ à $z_k$ dans $\handwritten{f}$, et en substituant $\handwritten{\unsure{\gamma|y}}_j$ à $u_j$ dans la série entière en $u_1,\ldots,u_q$ obtenue.

\uline{4}. \uline{La notion de fonction holomorphe}.

Le th.2 entraîne en particulier la proposition suivante :

PROPOSITION 4.- \uline{Soit} $\handwritten{f}(z_1,\ldots,z_p)=\handwritten{\sum}_{\handwritten{(n_k)}}\handwritten{\unsure{c|C}}_{n_1n_2\ldotp\ldotp n_p}z_1^{n_1}z_2^{n_2}\ldots z_p^{n_p}$ \uline{la somme d'une série entière sommable dans un polycylindre ouvert} $P$ \uline{de centre} $0$ \uline{défini par} $\handwritten{|}z_k\handwritten{|}\handwritten{<}r_k$ $(1\handwritten{\leq}k\handwritten{\leq}p)$. \uline{Soit} $(a_k)$ \uline{un point quelconque de} $P$ ; \uline{pour tout point} $(u_k)_{1\handwritten{\leq}k\handwritten{\leq}p}$ \uline{du polycylindre défini par} $\handwritten{|}u_k\handwritten{|}\handwritten{<}r_k-\handwritten{|}a_k\handwritten{|}$, \uline{les} $p$ \uline{polynomes} $a_k+u_k$ \uline{sont substituables aux} $z_k$ \uline{dans la série entière} $\handwritten{c}_{n_1\ldotp\ldotp n_p}z_1^{n_1}\ldotp\ldotp z_p^{n_p}$, \uline{et la somme de la série entière} (absolument sommable) \uline{par rapport aux} $u_k$ \uline{ainsi obtenue est égale à} $\handwritten{f}(a_1+u_1,\ldots,a_p+u_p)$.

En effet, avec les notations du th.2, on a ici $h_k(\handwritten{|}u_1\handwritten{|},\ldots,\handwritten{|}u_p\handwritten{|})==\handwritten{|}a_k\handwritten{|}+\handwritten{|}u_k\handwritten{|}\handwritten{<}r_k$ par hypothèse.

En d'autres termes, il existe un polycylindre ouvert de centre $(a_1,\ldots,a_p)$ dans lequel $\handwritten{f}(z_1,\ldots,z_p)$ est égale à la somme d'une série entière absolument sommable, par rapport à $z_k-a_k$ $(1\handwritten{\leq}k\handwritten{\leq}p)$.

DÉFINITION 2.- \uline{Soit} $A$ \uline{un ensemble ouvert dans} ${\unsure{K|\mathbb{K}}}^{p}$. \uline{On dit qu'une application} $\handwritten{f}$ \uline{de} $A$ \uline{dans} $\unsure{E|\mathbb{E}}$ \uline{est une fonction holomorphe dans} $A$ \uline{si}, \uline{pour tout point} $\handwritten{a}=(a_k)_{1\handwritten{\leq}k\handwritten{\leq}p}$ \uline{de} $A$, \uline{il existe un polycylindre ouvert de centre} $\handwritten{a}$, \uline{contenu dans} $A$, \uline{et dans lequel} $\handwritten{f}(z_1,z_2,\ldots,z_p)$ \uline{soit égale à la somme d'une série entière absolument sommable}, \uline{par rapport aux variables} $z_k-a_k$ $(1\handwritten{\leq}k\handwritten{\leq}p)$.

La prop.4 montre donc que :

PROPOSITION 5.- \uline{Si une série entière par rapport aux variables} $z_k$ $(1\handwritten{\leq}k\handwritten{\leq}p)$ \uline{est sommable dans un polycylindre ouvert de centre} $0$,
